\chapter{Speed-Qualified Consequences}

The running Alignment and its pose3 state now describe where the track frame
is and how it is oriented. They still do not determine a physical response.
Speed enters as a separate operating input; a vehicle model and evaluation
rule enter separately again.

\section{From Arc Length to Time}

For constant \(v>0\),
\[
\frac{\dd s}{\dd t}=v,
\qquad
\frac{\dd f}{\dd t}=v\frac{\dd f}{\dd s}.
\]
Thus the spatial curvature rate \(\kappa'(s)\) acquires a time rate only after
speed is supplied. The same holds for cant rate and higher derivatives.

\section{A Reduced Illustrative Relation}

For planar motion of a point following the reference trajectory,
\[
a_\kappa(s)=v^2\kappa(s).
\]
Under the local right-handed frame, small-scope rigid-body assumptions, the
cant convention of \cref{eq:derived-roll-angle}, and neglect of suspension,
wheel--rail, wind, and other effects, a reduced effective lateral component is
\[
a_{\mathrm{eff}}(s)=v^2\kappa(s)-g\sin\phi(s).
\]
This is an illustrative derived quantity, not a validated complete vehicle
dynamics model and not an authoritative cant-design rule.

\begin{table}[htbp]
\centering
\small
\begin{tabular}{lrr}
\toprule
Case at the same \(s\) & \(v^2\kappa\) & \(a_{\mathrm{eff}}\) for \(\phi=4^\circ\)\\
\midrule
\(v=20\,\mathrm{m/s}\), \(\kappa=0.001\,\mathrm{m^{-1}}\) & \(0.40\,\mathrm{m/s^2}\) & \(-0.28\,\mathrm{m/s^2}\)\\
\(v=40\,\mathrm{m/s}\), \(\kappa=0.001\,\mathrm{m^{-1}}\) & \(1.60\,\mathrm{m/s^2}\) & \(0.92\,\mathrm{m/s^2}\)\\
\bottomrule
\end{tabular}
\caption{Illustrative speed dependence for the same curvature and cant state (\(g=9.81\,\mathrm{m/s^2}\), rounded).}
\label{tab:same-curvature-two-speeds}
\end{table}

\Cref{tab:same-curvature-two-speeds} isolates the dependency: neither the
Alignment nor \(\kappa(s)\) changes, yet doubling speed multiplies the
curvature-induced term by four. Changing cant would change the effective
component and frame, not the horizontal curvature law.

\section{Variation and Operational Meaning}

With the same assumptions and constant speed,
\[
j(s)=v^3\kappa'(s)-gv\cos\phi(s)\,\phi'(s).
\]
This relation explains why transition shape and cant transition both matter
operationally while remaining distinct laws. If speed varies, additional
time-derivative terms appear. Vehicle response further depends on mass,
suspension, wheel--rail contact, and other model choices. None of those
quantities is silently folded into Alignment geometry.

\section{Evaluation Is a Later Operation}

The dependency order is
\begin{align*}
\text{geometry}&\rightarrow\text{realized frame}
\rightarrow\text{speed profile}\rightarrow\text{vehicle model}\\
&\rightarrow\text{derived dynamics}
\rightarrow\text{purpose and rule}\rightarrow\text{evaluation}.
\end{align*}
A limit from a standard is applicable only under its scope and version. A
computed exceedance is an evaluation result; acceptance, rejection, or
authorization requires an Engineering Decision with the relevant authority.

\begin{summary}
Speed changes the consequence of the running Alignment without changing its
intrinsic geometry. Cant changes the track frame and the reduced effective
action, while geometry, operating input, vehicle model, derivation, evaluation,
and decision remain separate.
\end{summary}
